\section{System Integration}\label{sec:integration}

Both joins run as one physical operator in Sirius, a GPU-native analytics engine with DuckDB~\cite{Yogatama26} as the frontend: DuckDB parses and optimizes a query, and Sirius executes the plan on the GPU.
Sirius supports datasets larger than GPU memory by partitioning tables into smaller batches which can automatically move across GPU memory, host memory, and disk.
This section describes how vector join is integrated into Sirius to support inputs larger than GPU memory (Section~\ref{subsec:integration-ooc}) and how vector join is optimized in the query plan (Section~\ref{subsec:integration-plan}).
% The disk tier is reached through the engine's scan and spill and is not evaluated in this paper.

\subsection{Inputs Larger than GPU Memory}\label{subsec:integration-ooc}

\yxy{This section is very hard to read. I thought it should be explaining how exact and approx joins can leverage Sirius to support data larger than memory. My understanding is this should be relatively easy since block NLJ and co-partitioned hash join can both support larger than memory originally. Maybe there is something else you are trying to discuss here?}
Both joins are out-of-core algorithms by construction: the blocked nested-loop join works on one pair of batches at a time, and the partitioned join on one partition at a time.
What Sirius adds is where the pieces live and how often they move.
The build side's codes or partitions stay on the GPU when they fit with a quarter of the engine's GPU pool to spare; otherwise they are kept in pinned host memory as ordinary engine batches, one per \emph{chunk} of rows, and each chunk a task needs is copied to the GPU, still coded, on its own stream so that the copy overlaps the previous chunk's work.
Nothing per build row stays on the GPU, and the fp32 rows of the candidates are gathered from the pinned column when they are re-scored.
A build side that was never loaded is read by the join's own scan, which codes each chunk to fp16 on the GPU as it arrives; the engine may spill scanned batches to disk and the join reads them back.
What the streaming costs is set by how many times the build side is read, so over a build side outside GPU memory the join inverts its loops: a task takes a \emph{probe block}, probe rows cut from several of the scan's batches and joined as one matrix, as many as a third of the GPU pool holds, and the build side is the inner input, read once per block, $\lceil m/b\rceil$ times for $m$ probe rows in blocks of $b$.
A self-join of one table scans it once and lets the probe side read the scan's batches.

\subsection{Planning the Join}\label{subsec:integration-plan}

In DuckDB's plan a vector join is a join: it has two inputs, a predicate, a cardinality estimate and a cost, so the optimizer's filter pushdown, join ordering and subquery handling apply to it as they apply to a hash join.
Sirius adds the decisions that have no relational counterpart.

\paragraph{Native SQL}
A vector join written in plain SQL reaches Sirius inside DuckDB's optimized plan, in one of the three shapes of Section~\ref{subsec:vector-join-interface}.
Sirius replaces the shape with the join operator and keeps the relational subtree that DuckDB built below the build side, such as filters, joins and subqueries.
It rewrites every expression above the join that reads a distance to read the join's score.
A plan that it cannot prove equivalent is left to DuckDB's CPU operators.
No data leaves the engine.

\paragraph{Access path}
Both decisions here start from DuckDB's cardinality estimate of the probe side, which sizes what the join builds.
An exact join with very few probe rows, too few to repay the code build, skips it and joins the fp32 rows directly; an approximate join takes its partition count from the estimate with the cost model of Section~\ref{subsec:approx-build}.
Codes and partitions stay for the session, so a wrong estimate costs at most one statement.


\paragraph{Join order}
A similarity predicate may sit beside a band predicate on a scalar column, such as two prices within 1\,\%.
DuckDB's own plan joins on the band and filters the pairs by similarity, which keeps the similarity predicate out of the vector join, and so off the GPU.
Sirius estimates the band's selectivity from column statistics and lets the band join go first only when it keeps so few pairs that DuckDB's cost for them is below the GEMM's; otherwise the vector join goes first and the band filters its output (Section~\ref{subsec:eval-planner}).
The decision is whether the similarity predicate becomes the vector join's own predicate, on the GPU, or a filter after a relational band join, on the CPU; it is not a choice between two GPU plans, and in our workload it has only ever fallen on the vector join's side.
A relational filter written outside the join's subquery is not pushed through the join and runs after it (Section~\ref{subsec:eval-limits}).

\paragraph{Aggregates inside the join}
A \texttt{count(*)} directly over a threshold join is computed in the kernel's epilogue, which counts the pairs within each row's bound and emits none.
A threshold join may also keep only each probe row's best match and report each pair of a self-join once, which a deduplication rule needs; the join then folds its candidates into one (distance, row) per probe row and re-scores only the candidates that can still win.
This form is reached through the table function today (Section~\ref{subsec:eval-limits}); a search API would need a top-2 search per row, since a row's nearest match in a self-join is itself, and a filter.
